\chapter{Leg Cramps (Nocturnal Leg Cramps)}

\section*{Definition}
Leg cramps are sudden, involuntary, and painful muscle contractions that most often affect the calf muscles but can also occur in the thighs or feet. They commonly happen at night (nocturnal leg cramps) and can disrupt sleep. The muscle feels hard and tender to the touch during a cramp.

\section*{What Happens in Your Body}
A muscle cramp occurs when a muscle contracts forcefully and does not relax. This is caused by abnormal nerve signals that overstimulate the muscle fibers. The exact mechanism is not fully understood, but it involves electrolyte imbalances, dehydration, or reduced blood flow affecting the nerve endings that control muscle contraction and relaxation. The cramp usually lasts from a few seconds to several minutes before the muscle finally relaxes.

\section*{Causes}
\begin{itemize}
  \item Idiopathic (no identifiable cause) — most common, especially at night
  \item Dehydration and electrolyte imbalances: Low potassium, magnesium, or calcium
  \item Muscle fatigue or overuse: Prolonged standing, walking, or exercise
  \item Sitting for long periods or poor posture
  \item Pregnancy (especially in the third trimester)
  \item Certain medications: Diuretics (water pills), statins, beta-agonists, and some blood pressure medications
  \item Medical conditions: Peripheral artery disease, diabetes, thyroid disorders, Parkinson's disease, cirrhosis, chronic kidney disease
  \item Alcohol use — can cause dehydration and electrolyte disturbances
  \item Vitamin D deficiency
  \item Nerve compression: Lumbar spinal stenosis or herniated disc
\end{itemize}

\section*{When to See a Doctor}
See your doctor if:
\begin{itemize}
  \item Cramps are severe, frequent, or interfere significantly with sleep
  \item You have swelling, redness, warmth, or tenderness in the leg — these could be signs of a blood clot (deep vein thrombosis)
  \item You have muscle weakness or loss of coordination in the leg
  \item Cramps are accompanied by leg pain that gets worse with walking and improves with rest (possible peripheral artery disease)
  \item The cramping does not improve with simple measures such as hydration and stretching
  \item You are taking a medication that could be causing the cramps
\end{itemize}

\section*{Related Symptoms}
\begin{itemize}
  \item Muscle soreness after the cramp subsides
  \item Visible muscle twitching or fasciculations
  \item Restless legs or leg discomfort at night
  \item Numbness or tingling in the feet
  \item Swelling in the legs or ankles
\end{itemize}
\textbf{Important:} If the cramp is accompanied by sudden swelling, redness, or warmth in one leg, seek medical attention — this could be a blood clot (DVT).

\section*{Self-Care / Home Management}
\begin{itemize}
  \item Stretch the affected muscle — for a calf cramp, straighten your leg and gently pull your toes back toward your shin.
  \item Walk around or jiggle the leg to help the muscle relax.
  \item Apply a warm towel or heating pad to relax tight muscles, or use an ice pack if the area is tender after the cramp.
  \item Massage the cramped muscle firmly but gently.
  \item Stay well-hydrated throughout the day, especially if you exercise or are in hot weather.
  \item Eat foods rich in potassium (bananas, potatoes, leafy greens), magnesium (nuts, seeds, whole grains), and calcium (dairy, fortified plant milks).
  \item Stretch your calf muscles before bed — stand at arm's length from a wall and lean forward keeping your heels on the ground, holding for 30 seconds.
  \item Keep blankets loose at the foot of the bed so your toes are not pushed downward.
\end{itemize}

\section*{How Your Doctor Will Evaluate You}
\begin{itemize}
  \item Your doctor will ask about the frequency, timing, and severity of cramps, as well as your activity level, fluid intake, and medications.
  \item They will check your legs for swelling, pulses, skin changes, and signs of nerve problems.
  \item Blood tests may be ordered to check electrolyte levels (potassium, magnesium, calcium), kidney function, and thyroid function.
  \item If peripheral artery disease is suspected, they may check your ankle-brachial index (comparing blood pressure in your ankle and arm).
  \item If a blood clot is suspected, an ultrasound of the leg veins will be ordered.
\end{itemize}

\section*{Treatment Options}
\begin{itemize}
  \item Most leg cramps resolve on their own without specific treatment and can be managed with stretching and hydration.
  \item Magnesium supplements may help some people, especially if deficiency is present — ask your doctor before starting.
  \item Quinine was previously used but is no longer recommended due to safety concerns (heart rhythm problems).
  \item Treating the underlying cause — correcting electrolyte imbalances, adjusting medications, or managing peripheral artery disease — usually resolves the cramps.
  \item Vitamin B complex supplements have shown some benefit in small studies.
  \item Stretching exercises, especially before bed, are the most effective prevention strategy.
\end{itemize}

\section*{References}

\begin{thebibliography}{99}
\bibitem{haldeman} Haldeman S, et al. Nocturnal leg cramps: a clinical review. \emph{JAMA}. 2023;329(18):1576-1585.
\bibitem{uptodate-cramps} Jansen PH, et al. Nocturnal leg cramps in adults. In: UpToDate, Post TW (Ed), UpToDate, Waltham, MA; 2025.
\bibitem{katzberg} Katzberg HD, et al. Muscle cramps: a practical approach. \emph{Muscle Nerve}. 2023;67(4):273-284.
\bibitem{allen} Allen RE, et al. Treatment of nocturnal leg cramps. \emph{BMJ}. 2024;384:e077538.
\bibitem{monderer} Monderer RS, et al. Sleep-related leg cramps: pathophysiology and management. \emph{Sleep Med Rev}. 2023;68:101749.
\bibitem{young} Young G, et al. \emph{Neurology in Clinical Practice}. 9th ed. Elsevier; 2023. Chapter 78: Muscle Cramps.
\end{thebibliography}
